
\documentclass{article}
\usepackage{colortbl}
\usepackage{makecell}
\usepackage{multirow}
\usepackage{supertabular}

\begin{document}

\newcounter{utterance}

\centering \large Interaction Transcript for game `cladder', experiment `full\_v1.5\_default', episode 7503 with Qwen3.6{-}27B.
\vspace{24pt}

{ \footnotesize  \setcounter{utterance}{1}
\setlength{\tabcolsep}{0pt}
\begin{supertabular}{c@{$\;$}|p{.15\linewidth}@{}p{.15\linewidth}p{.15\linewidth}p{.15\linewidth}p{.15\linewidth}p{.15\linewidth}}
    \# & \multicolumn{2}{c}{Player} && \multicolumn{2}{c}{Game Master} \\
    \hline

    \theutterance \stepcounter{utterance}  
    & & & \multicolumn{4}{p{0.6\linewidth}}{
        \cellcolor[rgb]{0.9,0.9,0.9}{
            \makecell[{{p{\linewidth}}}]{
                \texttt{\tiny{[P1$\langle$GM]}}
                \texttt{You are an expert in causal inference. The following question is not a typical commonsense query, but rather a meticulously designed question created by a professor specializing in causal inference, intended to assess the students' mastery of the course content.} \\
\texttt{Imagine a self{-}contained, hypothetical world with only the following conditions, and without any unmentioned factors or causal relationships: Wibl has a direct effect on vubr and zuph. Zuph has a direct effect on uvzi. Vubr has a direct effect on uvzi. For those who are not wibl and are not zuph, the probability of uvzi is 77\%. For those who are not wibl and are zuph, the probability of uvzi is 63\%. For those who are wibl and are not zuph, the probability of uvzi is 54\%. For those who are wibl and are zuph, the probability of uvzi is 76\%. The overall probability of wibl is 23\%. Will zuph increase the chance of uvzi?} \\
\texttt{Start your answer with 'Yes' or 'No', followed by additional reasoning or evidence to support your explanation.} \\
            }
        }
    }
     \\ \\

    \theutterance \stepcounter{utterance}  
    & \multicolumn{4}{p{0.6\linewidth}}{
        \cellcolor[rgb]{0.561,0.737,0.561}{
            \makecell[{{p{\linewidth}}}]{
                \texttt{\tiny{[P1$\rangle$GM]}}
                \texttt{No} \\
\\ 
\texttt{To determine whether `zuph` increases the chance of `uvzi`, we must calculate the causal effect of `zuph` on `uvzi`. In causal inference, this requires distinguishing between the **associational** effect (observed correlation) and the **causal** effect (interventional probability).} \\
\\ 
\texttt{**1. Analyze the Causal Structure**} \\
\texttt{The problem states:} \\
\texttt{*   `wibl` \$\textbackslash{}rightarrow\$ `vubr`} \\
\texttt{*   `wibl` \$\textbackslash{}rightarrow\$ `zuph`} \\
\texttt{*   `zuph` \$\textbackslash{}rightarrow\$ `uvzi`} \\
\texttt{*   `vubr` \$\textbackslash{}rightarrow\$ `uvzi`} \\
\\ 
\texttt{This structure implies that `wibl` is a common cause (confounder) of both `zuph` and `uvzi` (via the path `wibl` \$\textbackslash{}rightarrow\$ `vubr` \$\textbackslash{}rightarrow\$ `uvzi`). Therefore, `wibl` is a confounder for the relationship between `zuph` and `uvzi`. To find the causal effect of `zuph` on `uvzi`, we must adjust for (stratify by) `wibl`.} \\
\\ 
\texttt{**2. Calculate the Causal Effect (Average Treatment Effect)**} \\
\texttt{We need to compare the probability of `uvzi` when `zuph` is present versus when it is absent, holding `wibl` constant, and then average these effects over the population} \\
            }
        }
    }
    & & \\ \\

    \theutterance \stepcounter{utterance}  
    & & & \multicolumn{2}{p{0.3\linewidth}}{
        \cellcolor[rgb]{0.9,0.9,0.9}{
            \makecell[{{p{\linewidth}}}]{
                \texttt{\tiny{[GM$|$GM]}}
                \texttt{No} \\
            }
        }
    }
    & & \\ \\

    \theutterance \stepcounter{utterance}  
    & & & \multicolumn{2}{p{0.3\linewidth}}{
        \cellcolor[rgb]{0.9,0.9,0.9}{
            \makecell[{{p{\linewidth}}}]{
                \texttt{\tiny{[GM$|$GM]}}
                \texttt{no} \\
            }
        }
    }
    & & \\ \\

    \theutterance \stepcounter{utterance}  
    & & & \multicolumn{2}{p{0.3\linewidth}}{
        \cellcolor[rgb]{0.9,0.9,0.9}{
            \makecell[{{p{\linewidth}}}]{
                \texttt{\tiny{[GM$|$GM]}}
                \texttt{game\_result = WIN} \\
            }
        }
    }
    & & \\ \\

\end{supertabular}
}

\end{document}
